\documentclass[../../../include/open-logic-section]{subfiles}

\begin{document}

\olfileid[tr]{sth}{cardinals}{cardsasords}
\olsection{Sıra Sayıları Olarak Kardinal Sayılar}

Aslında kardinal sayılar teorimiz sıra sayıları teorimizi (utanmadan)
kullanacaktır. Yani kardinal sayıları yalnızca belirli sıra sayıları olarak
tanımlayacağız. Özellikle şu tanımı sunacağız:

\begin{defn}\ollabel{defcardinalasordinal}
$A$ iyi sıralanabiliyorsa $\card{A}$, $\cardeq{A}{\gamma}$ koşulunu
sağlayan en küçük $\gamma$ sıra sayısıdır. Her $\gamma$ sıra sayısı için
$\gamma=\card{\gamma}$ ise $\gamma$'ya bir \emph{kardinal sayı} deriz.
\end{defn}

Az önce ``$A$ iyi sıralanabilir'' ifadesini kullandık. Matematikte neredeyse
her zaman olduğu gibi buradaki kiplik ifadesi, bir varlık iddiasının
gayriresmî anlatımıdır: ``$A$ iyi sıralanabilir'' demek yalnızca ``$A$'yı
iyi sıralayan bir bağıntı vardır'' demektir.

Fakat \olref{defcardinalasordinal} tanımında bir engel vardır. \emph{Her}
kümenin bir büyüklüğü, yani her $A$ için $\card{A}$ bulunmasını isteriz.
Oysa verdiğimiz tanım bir koşulluyla başlar: ``$A$ iyi
sıralanabiliyorsa\ldots''. İyi sıralanamayan bir $A$ kümesi varsa tanımımız
$\card{A}$ nesnesini tanımlayamayacaktır.

Dolayısıyla \olref{defcardinalasordinal} tanımını kullanmak için her kümenin
iyi sıralanabileceğine ilişkin bir güvence gerekir. Ne yazık ki bu güvence
$\ZF$ içinde yoktur. Öyleyse \olref{defcardinalasordinal} tanımını
kullanmak istiyorsak şu türden yeni bir aksiyom eklemekten başka seçeneğimiz
yoktur:
\begin{axiom}[İyi Sıralama]
Her küme iyi sıralanabilir.
\end{axiom}
İyi Sıralama Aksiyomu'nun kabul edilebilir olup olmadığını
\olref[choice][]{chap} içinde tartışacağız. Fakat bundan sonra aksiyomu
serbestçe kullanacağız. Bununla kardinal sayıların
(\olref{defcardinalasordinal} içinde tanımlandığı anlamda) var olduğunu ve
iyi davrandığını ispatlamak oldukça kolaydır:

\begin{lem}\ollabel{lem:CardinalsExist}
Her $A$ kümesi için:
\begin{enumerate}
	\item\ollabel{cardaexists} $\card{A}$ vardır ve biriciktir;
	\item\ollabel{cardaapprox} $\cardeq{\card{A}}{A}$;
	\item\ollabel{cardaidem} $\card{A}$ bir kardinal sayıdır; yani
	$\card{A}=\card{\card{A}}$.
\end{enumerate}
\end{lem}

\begin{proof}
$A$ sabit olsun. İyi Sıralama gereği bir $\tuple{A,R}$ iyi sıralaması
vardır. \olref[ordinals][ordtype]{thmOrdinalRepresentation} gereği
$\tuple{A,R}$ biricik bir $\beta$ sıra sayısıyla izomorftur. Dolayısıyla
$\cardeq{A}{\beta}$ olur. Sonlu Ötesi Tümevarım gereği
$\cardeq{A}{\gamma}$ olan biricik en küçük $\gamma$ sıra sayısı vardır.
Öyleyse $\card{A}=\gamma$ olur ve \olref{cardaexists} ile
\olref{cardaapprox} kurulur. \olref{cardaidem} için $\delta\in\gamma$ ise
$\gamma$'nın seçimi gereği $\cardless{\delta}{A}$ olduğunu gözlemleyin.
Eş güçlülük bir denklik bağıntısı olduğundan
(\olref[sfr][siz][equ]{equinumerosityisequi})
$\cardless{\delta}{\gamma}$ da olur. Dolayısıyla
$\gamma=\card{\gamma}$.
%So, for reductio, suppose that there is some ordinal $\delta \in \gamma$ such that $\cardeq{\gamma}{\delta}$. Then, $\cardeq{\cardeq{A}{\gamma}}{\delta}$ so that $\cardeq{A}{\delta}$ by \olref[sfr][set][equ]{equinumerosityisequi}, which contradicts the choice of $\gamma$ as the least ordinal such that $\cardeq{A}{\gamma}$. 
\end{proof}

Sonraki sonuç Cantor İlkesi'ni ve daha fazlasını güvence altına alır.
(Kardinal sayıların sıralamalarını sıra sayılarından devraldığına dikkat
edin; yani $\cardfont{a}<\cardfont{b}$ ancak ve ancak
$\cardfont{a}\in\cardfont{b}$ olur. Bunu formüle ederken kardinal sayı
olduğunu bildiğimiz nesneler için Fraktur harfleri kullanacağız. Bu oldukça
standarttır. Yaygın bir alternatif, kardinal sayılar sıra sayıları olduğundan
Yunan alfabesinin ortalarından $\kappa,\lambda$ gibi harfler seçmektir.):
\begin{lem}\ollabel{lem:CardinalsBehaveRight}
Her $A$ ve $B$ kümesi için:
\begin{align*}
	\cardeq{A}{B} &\text{ ancak ve ancak } \card{A} = \card{B}\\
	\cardle{A}{B} &\text{ ancak ve ancak } \card{A} \leq \card{B}\\
	\cardless{A}{B}&\text{ ancak ve ancak } \card{A} < \card{B}.
\end{align*}
\end{lem}

\begin{proof}
İkinci iddianın soldan sağa yönünü ispatlayacağız (öteki durumlar benzerdir
ve alıştırma olarak bırakılmıştır). Şu diyagramı ele alın:
\begin{center}
	\begin{tikzpicture}
	\node (nodea) {$A$};
	\node[right = 6em of nodea] (nodeb) {$B$};
	\node[below = 2em of nodea] (nodecarda) {$\card{A}$};
	\node[below = 2em of nodeb] (nodecardb) {$\card{B}$};
	\draw[->] (nodea)--(nodeb);
	\draw[<->] (nodea)--(nodecarda);
	\draw[<->] (nodeb)--(nodecardb);
	\draw[->, dashed] (nodecarda)--(nodecardb);
\end{tikzpicture}
\end{center}
Çift başlı oklar, varlıkları \olref{lem:CardinalsExist} tarafından güvence
altına alınan !!{bijection}s gösterir. $\cardle{A}{B}$ varsayımı altında
bir !!a{injection} $A\to B$ vardır. Okları $\card{A}$'dan $A$'ya,
$B$'ye ve $\card{B}$'ye kadar izlediğimizde bir !!a{injection}
$\card{A}\to\card{B}$ (kesikli ok) elde ederiz.
\end{proof}\noindent
\olref{lem:CardinalsBehaveRight} sonucunu Schröder--Bernstein'ı yeniden
ispatlamak için de kullanabiliriz. İddia, $\cardle{A}{B}$ ve
$\cardle{B}{A}$ ise $\cardeq{A}{B}$ olduğudur. Bunu
\olref[sfr][siz][sb]{thm:schroder-bernstein} olarak ifade etmiş, fakat ilk
kez \olref[sfr][infinite][card-sb]{sec} içinde biraz çabayla
ispatlamıştık. Şimdi şuna bakalım:

\begin{proof}[Schröder--Bernstein'ın Yeniden İspatı]
$\cardle{A}{B}$ ve $\cardle{B}{A}$ ise
\olref{lem:CardinalsBehaveRight} gereği $\card{A}\leq\card{B}$ ve
$\card{B}\leq\card{A}$ olur. Üçlü Karşılaştırma ve
\olref{lem:CardinalsBehaveRight} gereği $\card{A}=\card{B}$ ve
$\cardeq{A}{B}$ elde edilir.
\end{proof}
\noindent
Bu çok basit bir ispat olsa da örtük olarak hem Yerine Koyma'ya
(\olref[ordinals][ordtype]{thmOrdinalRepresentation} sonucunu elde etmek
için) hem İyi Sıralama'ya (\olref{lem:CardinalsBehaveRight} sonucunu
güvence altına almak için) dayanır. Buna karşılık
\olref[sfr][infinite][card-sb]{sec} ispatı çok daha bağımsızdır (hatta
$\Zminus$ içinde yürütülebilir).

\end{document}
